
% 来源：EPUB part124.xhtml
\section{苹果公司}

\subsection{喜欢苹果的理由}

说说我喜欢苹果的一些理由。这不是论文，想到就说，没有重点和先后次序。

1. 苹果的产品确实把用户体验或消费者导向做到极致了，对手在相当长的时间里难以超越甚至接近（对喜欢苹果的用户而言）。

2. 苹果的平台建立起来了，或者说生意模式或者说护城河已经形成了（光软件一年都几十亿的收入了）。

3. 苹果单一产品的模式实际上是我们这个行业里的最高境界，以前我大概只见到任天堂做到过（索尼的游戏产品类似）。单一产品的模式有非常多的好处：

（1）可以集中人力物力将产品做得更好。比较一下iPhone系列和诺基亚系列（今年要推出40个品种）。苹果产品的单位开发成本是非常低的，但单个产品的开发费却是最高的。

（2）材料成本低且质量好，大规模带来的效益。苹果的成本控制也是做到极致的，同样功能的硬件恐怕没人能达到苹果的成本。

（3）渠道成本低。呵呵，不是同行的不一定能明白这话到底有什么分量（同行也未必明白），我是20年前从任天堂那里学会的。那时很多做游戏机的都喜欢做很多品种，最后下场都不太好。

4. 苹果的营销也是做到极致了，连广告费都比同行低很多，卖的价钱却往往很好。

5. 苹果的产品处在一个巨大并还有巨大成长的市场里。

（1）智能手机市场有多大？你懂的！

（2）平板市场有多大？你也会懂的。

总而言之，我认为苹果现在其实还处在其成长的早期，应该还有很大的空间。（2011-01-22）

实际上，排名不是我们要追求的东西。能做出用户喜欢的产品才是我们的追求。也许真有哪一天我们也能做出个伟大的产品，希望50年内。（2010-06-08）

不在乎排名的意思是我们只在乎用户的感受，排名只是结果而已，不是追求的目标。（2010-06-09）

上面两句话是在买苹果之前说的。苹果确实做到了我想做的事情，2011年1月的时候我就是想通了这点才开始买苹果的。（2024-09-12）

说一点我特别喜欢苹果的地方：苹果是非常难见到的能够长期聚焦在做对的事情上的企业。

多数人在投资时很习惯去看有没有把事情做对这点，从而会很容易掉进短期表现当中。即使是伟大的企业，在做对的事情的过程当中，也是可能会犯错的，因为把事情做对往往需要一个过程。所谓长期投资，从某种角度上说，就是要能够相信坚持做对的事情的、有能力的公司，最后会把事情做对。（2013-02-13）

OPPO和苹果其实有很多相同的基因，这也是我最后能看懂苹果的原因之一。（2011-08-07）

\BookQuestion{网友：巴菲特卖IBM，买入苹果，到底是看到了什么本质？}

现金流，现在的现金流和未来的现金流。（2018-03-04）

% 来源：EPUB part125.xhtml
\subsection{苹果最厉害的是生态系统}

苹果最厉害的就是在苹果的企业文化下建立的生态系统，这是个非常强大的商业模式，非常难以撼动。（2020-11-16）

\BookQuestion{网友：苹果公司在乔布斯回归之后，推出iPhone之前，是不是也算企业文化强大的苦生意？}

至少iPhone出来后苹果的商业模式变得更好了，生态圈建成了。（2025-01-25）

有些东西进步到一定程度后其实后面的变化就不大了，比如20世纪30到40年代的汽车发动机和现在其实没本质差别。有些国家（虽然我们很多人很看不上）40年代就能造航空母舰，说明由于某些原因，时间有时候帮不上大忙。

苹果最厉害的是现在这个平台以及生意模式，以我的水平我看不出谁能打破它。当然有很多以为水平比我高的人非说别人可以，当然也有自己说自己可以打败苹果的，你看看他们像吗？

也许，能打败苹果的人快出生了，也许正在某个地方徘徊，但怎么看都不是现在大家知道的这些，这是个大家可以看很多年的故事。（2012-01-28）

苹果厉害的不仅仅是硬件，这是所有其他竞争对手没辙的地方。三星顶多不过是新的诺基亚而已，对苹果没有威胁。苹果的生意模式实在是强大，其他的“对手”已经难以撼动苹果了。（2012-02-05）

除苹果以外的卖智能手机的最后结局可能会有点像卖PC的，多数大概都生存不下来。（2012-03-08）

其实苹果的强大实力是不容易也不应该被忽视的，除非你有意去忽视，但鸵鸟的办法不会有任何帮助。任何一个被苹果直接或间接、现在或未来影响到的行业里的企业，大概早晚都要面对苹果的强大实力。随便举几个已经受到影响的行业或产品吧：手机、电脑、手提电脑、游戏（包括网游）、游戏机、中低端相机和录像机、音乐、电影、书、电视、网络、运营商、所有相关产品的零售、几乎所有的电子消费品……欢迎补充。找到在苹果生态系统里生存的方式也许比挑战苹果要来得更现实一些。能够避开苹果也许是另一个办法。（2012-03-12）

\BookQuestion{网友：如果苹果某次发布的新产品失败了，那对苹果的影响应该很大，这时候它还是好生意吗？虽然苹果现在的客户黏性很大，但是我想也经不起连续2年的新产品失败吧？}

你见过在别的同行不停推新品的同时能够按兵不动一年营业额依然不掉的公司吗？而且这还是在上个产品不那么轰动的前提下。如果这还不叫模式好，那什么叫模式好？现在轮到苹果出招了。顺便补充一下，凡是认为iPhone5是失败产品的人都应该问一下在人类历史上有过几个单一产品比iPhone5成功的？（2013-09-04）

\BookQuestion{网友：用习惯iPhone的人，一般改用安卓系统会不习惯的。苹果只要一输入ID，程序、音乐、通讯录等等都会出现，通过AppleTV，可以把其他设备上的图像、视频、音乐随时显示出来。通过查找iPhone随时可以知道在外地的亲人在哪，然后通过FaceTime免费视频。当iMac电脑的图像一键显示到投影上，当随时在中国用iPhone启动在美国家中的电脑并打开摄像头查看自己的宠物有没有捣乱，你就会感觉到科技的力量。苹果成为世界第一大公司是理所当然的！还有苹果绝对是消费者导向的，你想要的功能它基本都已经实现了。安卓系统这些功能可能借助其他应用也可以实现，但感觉一切都乱，设备互联性太差。苹果是个大系统，音乐、视频、游戏、软件、相机、电脑等的钱它全赚了，关键你买的音乐、视频、程序等，一辈子都是你的，但必须在它的设备上才可以用。谁还会轻易换设备？AppleTV一定要买，它可以把所有苹果的设备都联起来。}

你说得很清楚，这个就叫生态系统。安卓的问题是现在卖得越多，将来可能就越乱，慢慢跑到苹果生态系统里的人就会越多。

还有一个重要的东西，就是苹果的安全性。那么多企业要用iPhone和iPad是有道理的。美国500强里95\%开始用iPhone和iPad系统，世界500强里85\%开始用——库克在股东大会上讲的。

听说目前好像只有iPhone不能被装窃听器，其他的手机绝大部分好像都可以，怕被人窃听的要小心了。（2013-03-04）

\BookQuestion{网友：国内支持Apple Pay的地方不多，不过也不用带信用卡，零钱也不用带，微信支付能覆盖所有场景。}

我有微信支付，也知道微信支付确实很成功，不然我这腾讯股东就白当了。不过，Apple Pay在很多国家里已经慢慢普及开了，国内慢慢也会有越来越多人用的，但因为有微信支付等的原因，估计要慢很多，但10年后也许也会很普及的。我想说的是，苹果系统未来在很多国家和地方会在金融体系里起到很大作用，当然这个长时间内不会在国内体现的。（2022-04-02）

\BookQuestion{网友：你对你们公司怎么定位？}

有一天我们也有可能在这个行业生存不下去的，安卓手机的差异化未来会越来越难。当然，我们也有可能会成为生存下来的不多的几家之一。

\BookQuestion{网友：现在华为也大举推广智能机了，是不是会挤压其他厂商？}

彼此挤压吧。这也是导致最后大多数人都没钱赚的原因之一。（2012-03-10）

\BookQuestion{网友：为什么安卓的体验总是差强人意？}

安卓本身占用资源大，比如如果都是两核或四核的手机，表现就要差不少，加上大家看到的苹果的东西是几十年积累下来的，安卓要全面达到恐怕不太可能。但是，安卓的表现其实已经很不错了，已经比当年诺基亚智能机的感觉强了许多倍，某些使用方面和苹果相比各有特色也是一定会有的，而且价钱还可以便宜很多。未来安卓手机之间的差异化会越来越小，三星的价钱要挺住不容易。苹果的生态系统会保住特别喜欢或习惯的用户，所以受安卓的影响不会像想象的那么大。不过，有苹果和三星在上面顶着价格，其他卖安卓智能机的企业应该还有一段不错的小日子。（2013-01-23）

\BookQuestion{网友：您如何看三星的生意模式？}

三星的生意模式其实不太好，属于疲于奔命的那种。三星的问题在于其实没有产品追求，是生意导向的公司，焦点太发散，时间长了大概率会出问题。不过韩国企业非常努力，会把危机延长到比较大才出来，这点和日本企业比较像。

看三星和苹果有点像看篮球里美国队对其他队的比赛，场面上很激烈，其实差距很大。（2013-02-13）

\BookQuestion{网友：目前智能机市场，一种模式是类似三星这种“多机型”模式，另外一种是类似苹果的“精品”模式。哪一种模式更好，为什么？}

有本事做成苹果的会愿意做三星吗？（2013-08-02）

苹果倒是像卖车的，一年出一款。三星像诺基亚，结局可能也像，不信过10到15年再看。（2012-06-30）

\BookQuestion{网友：1. 三星、HTC一直搞机海战术，其实大部分厂商都是这招，苹果出来后，这招不灵了。HTC今年出one系列，说不玩换壳了，真能转性吗？2. 苹果真是独步江湖了。安卓阵营跟苹果的差距，不是简单的机子外表问题，是整个商业模式和生态系统的问题。如果专注研发ROM（Read- Only Memory，只读存储器），把消费者体验提升有戏吗？3. 我发现国内很多用iPhone的只是觉得很漂亮或者纯粹是炫耀性消费，iPhone的威力他们恐怕一小半都没用到。}

1. 人大概是不会变的。

2. 苹果以外的市场依然很大，而且上升空间也很大。专注于用户体验的企业总是会走得远一些的。

3. 我也发现很多人还没有发现iPhone的很多好处就已经很喜欢iPhone了，所以觉得iPhone的潜力还很大。

\BookQuestion{网友：您评估苹果的时候，有没有担心有一天会有新的颠覆性的技术出现，取代现有的东西？}

没有天上掉下来的技术。如果真的有重大的技术变革，我相信我是能够看得见的，而且我认为这种变化会大概率地来自苹果。（2014-02-06）

\BookQuestion{网友：以后会不会有新技术代替手机？}

一定会的，但不会很快。你需要想明白的是谁，什么时候，等等。（2015-11-10）

\BookQuestion{网友：有可能是苹果也不出奇。}

以前我们做VCD时，有人向我挑战说DVD将来会取代VCD，我说你觉得将来卖DVD的会是谁？（2015-11-11）

\BookQuestion{网友：很多公司市值超过5000亿美元之后就会增长缓慢或走向衰落，苹果是不是也一样？苹果已经是地球最大公司，不论它现在基因多么优秀，最终会因庞大而死亡。这个公司已经没有增长空间了，除非地球经济飞速增长，否则无法支持这样的市值了。}

科技类活不长指的是微软吗？即使今天来看，如果从某一时点用同样代价让你可以拥有1\%的微软或可乐的股份，哪个会让你觉得更开心呢？现在很多科技类的公司走下坡路实际上都是同一个原因，你懂的。还有就是，消费类的公司就一定活得长吗？只要看看可口可乐这么多年的同行还有多少活下来就明白了。（2012-04-29）

\BookQuestion{网友：苹果是否会盛极而衰？}

多年来人们一直这么说，可他们就是不明白。（2012-03-29）

% 来源：EPUB part126.xhtml
\subsection{做最好的产品}

昨天去苹果的股东大会，听到库克讲了类似的话，大意是：虽然我们可以有个按钮，可以随时做最多的产品，但我们的文化是做最好的产品而不是最多的。不过，他也顺便提了一下，认为要建立一个强大的生态系统，必须要有一定的量。所以苹果是会兼顾最好的产品和足够的份额之间的关系的。

另外，库克也提到量的概念其实不仅仅是卖了多少，而是人们买了以后用多少。有兴趣的可以查下新浪上微博的活跃用户都在用啥，也许大家会发现用苹果的人的比例其实远远超过印象中的市场占有率。

参考的统计办法：可以看热门微博Top100里发帖人是用什么发的帖，还可以看评论人用什么写的评论，也可以看微话题下面评论人是用什么评论的。

\BookQuestion{网友：刚才数了一下，Top100微博：来自iPhone20个；来自iPad5个；来自Android3个；来自“其他手机”1个；其他的是显示来自微博，或某个应用某个网站之类。}

看一下娱乐版的会更有趣些。这个比例说明苹果产品的使用率明显高于其他产品，好像没有别的指标可以更好地鉴别一个产品的好坏了。在微话题下，苹果的使用量大概接近50\%，其他的所有移动设备占到50\%多一点，这些都是普通上网者。

记得库克还说：我们最重要的是聚焦在做最好的产品上，利润会随之而来的。（2013-02-23）

\BookQuestion{网友：当年做PC时，封闭的苹果输给开放的微软 + IBM，今天的苹果一样封闭，难道它一定能赢谷歌 + 三星吗？}

我觉得不一样。你用一用产品就明白了，不是一回事。（2011-01-23）

我不是所谓的“果粉”，像我这个年纪和这么理性的人，不大会是任何人或东西的粉丝，我甚至不是巴菲特的粉丝。实际上我以前一直多少有点不太喜欢苹果，毕竟用了PC和诺基亚这么多年，但iPhone和iPad的出现确实彻底改变了我对苹果的认识。作为一个普通人，我没能力从技术上去说这些产品，所以只说说我都用iPhone和iPad干什么。

1. 首先，iPhone和iPad都非常容易上手。早年用摩托罗拉的时候，心里那个烦啊，用户界面太难用了。后来用诺基亚的智能手机时又找到了当年用摩托罗拉的感觉。以我个人的水平而言，安卓的界面水平还远远比不上iPhone（也许我们自己的出来后会改变一点这个局面？期待中）。上个月我曾经想买个安卓手机试试的，在商店里转了半个多小时，店员给我介绍了10来分钟，在3个牌子（HTC、三星、摩托罗拉）间挑来挑去，晕晕乎乎中最后还是决定不给自己添麻烦了。

iPhone？呵呵，我几乎拿到手里时就会用了。店员只是告诉我开关在哪里，那些图标是什么，前后花了大概15秒钟。当苹果的销售员可真容易啊，我自己前段时间在苹果店里只花了几分钟就推销了一台iPad2。

2. 我用iPhone干什么？我的iPhone上下载的软件其实比大多数人都少，毕竟不太花时间玩这个东西。电话、短信、E-mail、上网、世界时间、天气、图片、地图、记事本、YouTube、相机、录音、录像、字典、相册、图书馆、iPod、FaceTime、电影、电视、股票查询及交易、报纸、指南针、日历（记事本）、游戏，等等，不太注意但特别常用的是手电筒。这些功能可能只是iPhone的沧海一粟。其实我想说的不是iPhone有多少功能，最重要的是所有这些功能都非常容易上手。一学就会，无师自通！

3. 很多人说iPad其实就是个大号的iPhone。在用iPad之前我同意这个说法，但用了一段时间后觉得这个大确实大得非常有道理。我现在基本上已经不太用我的手提电脑了，旅行的时候带上iPad就绝对够了，iPhone有时候还是觉得有点小。对我而言，iPad2最强大的地方就是网络功能很强大，上网非常快，又没有手提电脑开机的麻烦，用起来真是很爽，携带起来又轻便。还有一个特别强大的功能就是图片，包括上网看图片。我的iPad2里有大概1万张照片，记录着我们生活中的很多瞬间，没事看看真的觉得生活很美好，这个玩意儿该值多少钱呢？FaceTime也是一个很好的功能，在有Wi- Fi的前提下用起来那个叫方便，比打电话还容易。

我给我妈妈也买了个iPad2，最近能让人带回去，里面也装上我们的照片。以后视频通话就非常方便了（以前曾经想过用电脑视频，实在是有点麻烦），没事老人家还可以翻翻照片，真是好啊。

别的就不说了，有机会大家自己体会吧。总的来讲，如果经济上不是个太大的负担的话，iPhone和iPad2是绝对值得拥有的产品，生活确实因此变得不同了。（2011-03-31）

不知道大家有没有丢电话的经历？丢了以后怎么办？

前段时间我在旅行中丢了一件行李（iPad在里面），当时以为是丢在机场了，后来突然想起苹果的Find my iPhone的功能，结果用iPhone一查（可以用任何别的IOS设备），发现我的iPad还在酒店里。（值钱吧？）马上打电话给酒店告诉了他们，结果几天后酒店就把行李寄回来了。还有一次把iPhone掉在沙发的缝里了，一时半会找不到，结果还是靠Find my iPhone找到的。（这个功能你愿意付多少？）

当然，就算找到在哪里，万一捡到的人不肯归还怎么办？Find my iPhone上还有一个功能，那就是可以立刻消除你丢失的iPhone或iPad上的所有资料。另外就是，在补一个新的iPhone或iPad时，只要用自己的AppleID，几乎所有的资料就又都回来了。（这个功能值多少钱？）

有个朋友前几天说丢了电话所以丢了所有电话号码，对于iPhone使用者来说，这个事情不会再发生了。（这个值钱不？）

好像只要记得自己的AppleID就可以用这个办法找到自己的iPhone或iPad，只要开始用的时候把ID开通了就行。（这个其实也很有价值，因为人们往往是丢了以后才知道这个功能的。）

都说iPhone贵，想想丢电话时自己愿意付多少钱把丢掉的资料找回来并且不失密吧。如果觉得这个重要的话，那就……iCloud确实强大，不然三星的人不会回家后用iPhone和iPad的。（2012-12-13）

\BookQuestion{网友：哪位朋友能谈谈iWatch？}

那还是一个传说，从何谈起？不过，关于苹果的传说最搞笑的就是廉价iPhone。有了这个传说后，接着就有了苹果该不该这样做的争论。其实讨论苹果该不该卖廉价iPhone的人大概都是不太懂苹果的，因为苹果的文化里根本就不会有这个争论。（2013-02-15）

\BookQuestion{网友：关于苹果AR眼镜和汽车有什么可以分享的么？}

我不知道苹果在搞什么，但确实相信苹果产品最终都会是精品。（2020-12-06）

Vision Pro对喜欢它的用户而言是个“nice to have”的产品，而iPhone则是“need have”。Vision Pro要上升到“need have”我猜可能需要10年。

我说10年的意思是需要很长的时间，就是短时间里概率很低的意思。（2024-06-29）

\BookQuestion{网友：我这样理解苹果，现在苹果已经不是靠一种新产品就可以打败的了，因为它有良好的企业文化（用户体验）、生态链（App应用商店、电脑和手机、平板共享平台），等等。}

其实革命性的产品并不会贴个标签在上面。iPad mini绝对是个革命性的产品，大家会看到的。另外，苹果的产品是非常细致的。三星最近据说要推8寸的平板，比苹果7.9寸的大一点。别的不说，有一点可能就会很有趣。iPad mini我是可以放进裤子口袋的，但只是刚刚好放进去，个人觉得iPad mini可能是测过的，8寸的很可能放不下。（2013-03-01）

\BookQuestion{网友：从苹果的哪些做法中能看出下一年推大屏的可能性很大？}

我知道大屏市场有多大，苹果当然也知道！只要在苹果的原则范围内和能力圈内，苹果没道理不推。不过，我可能快两年前就说过苹果必推大屏，没想到等了这么久。5s出来后，苹果大概就可以抽出手来推大屏了，所以我认为苹果一年内就会推的。（2013-09-16）

\BookQuestion{网友：乔布斯说的Think different是什么意思？}

Think different是不从众的意思。这个说法并不完全对。正确的说法是消费者导向。当年苹果一直不出大屏手机就是个例子。消费者需要大屏，但苹果Think different，结果晚了3年多。（2020-11-26）

其实大屏手机也从另外一个角度说明了苹果厉害的地方。可能因为乔布斯说过不要大屏，所以苹果犹豫了一下，但消费者的需求摆在那里，所以最后他们还是用不同的办法想通了。我和方三文曾经打过一个小赌（没赌任何东西）。方三文问：你觉得苹果一定会出大屏手机吗？我说：他们一定会的！果然3年后他们就出了。苹果大屏手机出得晚给了我很多时间去思考苹果和买苹果，不然我可能要少买很多苹果哈。（2020-11-27）

\BookQuestion{网友：今天看到苹果官方出售二手iPhone的消息，包括各型号的回收价与售价对比，段总怎么看？}

我的看法是，只有苹果能够做这件事情，别人其实做不了。（2015-04-14）

苹果产品的残值高也说明苹果的东西确实好。（2013-12-13）

\BookQuestion{网友：大道怎么看Meta Orion，会是手机之后的下一代爆品吗？}

他们完全不是苹果的对手。我用过他们以前的产品，确实不是一个级别的。（2024-09-28）

\BookQuestion{网友：您对安卓的前景怎么看呢？安卓做高端有前途吗？另外您觉得苹果会不会着力占领廉价手机市场？这样做的话对他们好处多还是坏处多？}

在相当长的时间里，安卓的量依然会是大的。不能说安卓做高端如何，但某个品牌用安卓做个很贵的手机出来是很可能的，会买的人当然要叫那个高端了。另外，在安卓里也有高低端之分。

苹果大概永远都不会专门做所谓的廉价手机市场，对你提出这个问题感到有点意外，说明你对企业的理解已经和大众很接近了。苹果的策略一直都在那里摆着呢，做最好的产品，让前面的产品做入门级以及提高级的产品。大约3年内，当iPhone7出来的时候，iPhone5c就会变成入门级的产品，大概会卖到2000出头，那个时候会有点好看哈。

另外，iPhone5s绝对是划时代的，不知道为什么那么多人会视而不见？我可以肯定的是，三星里很多人要开始“熬夜多压力大”了。（2013-09-15）

% 来源：EPUB part127.xhtml
\subsection{iPhone很可能是最便宜的手机}

我一直觉得iPhone很可能是最便宜的手机，因为单位时间里（比如10年20年）需要花在买手机上的钱的总额很可能是最少的，如果你想省钱的话。（2023-11-03）

\BookQuestion{网友：成为街机的iPhone的热度会持续多久？已经无法成为身份象征的iPhone是否还会得到富裕人群的青睐？}

一个季度3700万富人买iPhone？买iPhone只是因为iPhone好用，只有不用iPhone的人才可能认为iPhone是身份象征。（2012-01-31）

\BookQuestion{网友：我是果粉，家里用苹果和其他品牌的产品。就产品而言，无论技术、外观、感觉和吸引的手段，苹果打100分，其他品牌打85分。苹果公司在品牌上有很强的优势，很高的忠诚度，这些是其他品牌难以追赶的。}

对于价格不那么敏感的用户而言，谁愿意买85分的东西而不是100分的东西，尤其是每天都要拿在手上用的手机呢？（2012-02-03）

\BookQuestion{网友：iPad、iPhone、iPad mini是我拥有的最超值的产品，喜欢用的功能都做到了让我非常满足的程度。要是没有苹果，我花更多的钱和时间也达不到它们带来的效果。这两年很轻松地就推荐出去十几部，还没有一个抱怨买亏了的。}

是哦，看起来免费的东西往往最贵，看起来便宜的东西往往不便宜。买东西如此，投资也常常如此。我也觉得i系列是最便宜的东西，一个是消费体验好（这个值很多钱哈），换机的钱也省了。（2013-02-28）

\BookQuestion{网友：免费的东西往往不便宜，贵的东西有时其实很便宜。比如，买本好书才几十块钱，收获可能远超几百倍书的价钱。再比如您请巴菲特吃饭，巴菲特说最重要的是生意模式，您说收获远超饭局的价钱。我现在也更看得懂了，既是很好的投资，又顺带做了慈善。}

能明白贵与便宜的关系非常不容易啊，绝大部分人大概一辈子都很难明白的。（2013-03-29）

特斯拉看起来很酷，但有点贵，不是那么容易买得起。iPhone虽然相对贵不少，但绝对值差别非常小。（2013-05-14）

\BookQuestion{网友：看不懂“相对贵”和“绝对值”？}

你两个都多用一段时间就明白了。1000块和800块的绝对值差异是200块。如果1000块的东西好用那么一点点，带来的回报会远超过200块。大多数人会觉得800块的东西便宜，实际上往往相反。（2013-05-23）

\BookQuestion{网友：我的观察，买苹果的人大部分不是因为性价比高。人们用苹果主要是因为好用、放心用、有面子。性价比其次。当然，有很多人因为售价高而不舍得买的。但这部分人也不是因为性价比不够高而不买的。}

我同意苹果产品是高性价比的。高性价比的东西不等于便宜，该知道的人最后都会知道的。好的车往往也是性价比高的车，而且这些车往往比较贵，想通这个就明白苹果的性价比问题了。当然，对于那些整天换手机，最后换到iPhone的人们而言，这个问题就不是问题。（2014-08-28）

每个人都是会看性价比的，但每个人对价格的敏感度不同。还有就是每个人对钱的看法不一样。比如有人会为了便宜去买不安全的车，等等。（2013-02-27）

% 来源：EPUB part128.xhtml
\subsection{单一产品模式}

单一比较容易聚焦，这意味着把事情做对的概率高同时犯错误的概率低。

要做到单一是非常难的一件事，因为市场的需求是多样化的。iPad和iPhone有一样神奇的东西，那就是不管什么年龄段的用户都觉得这是个专门为自己设计的产品。

我那个“著名”朋友李录是个孝子，给70多的老爸买了个iPad，结果老爸一天到晚地捧着iPad，还跟李录说：这东东简直就是专门为我设计的嘛。（2011-04-09）

\BookQuestion{网友：当年的小霸王学习机也差不多是单一产品模式，记得当时单一款486B卖得很好，而且在旺季来临前就备足了货源。}

呵呵，这是我能看懂苹果的重要原因之一，没有那段经历，不一定会知道苹果有多厉害。（2012-02-01）

\BookQuestion{网友：为什么单一产品模式可以降低渠道成本？}

你倒过来想也许就明白了。举个例子，20多年前我们卖小霸王游戏机（还没到学习机）时，小霸王游戏机只有两款：红白机和蓝白机（有点像iPhone的黑色和白色），对手们经常都是有好几十款的。当时我们的销售部门被质疑最多的问题就是款式太少。好像最后款式多的那些对手绝大多数都消失了，原因大概就是品种太多最后导致的问题（其实是表象）。（2013-02-16）

这里说的渠道成本的意思就是品种越多，占用的渠道库存越大，是几何级数的关系。（2011-04-09）

一般来说，因为技术进步非常快，电子产品一旦一定时间内卖不出去贬值就会非常高。同一品牌的产品的品类在超过一定数量后，库存会比较难管理，处理成本会非常高。（2020-10-29）

\BookQuestion{网友：品种越多，占用的渠道库存越大，是不是说会给经销商增加负担？}

经销商的负担难道不是自己的负担？怪不得有人为了2015年成第一而向经销商压货呢。（2011-04-09）

\BookQuestion{网友：您当年是怎么控制库存的？}

安全第一的原则。小步快跑，宁愿少赚些，不要出大问题。（2011-12-15）

\BookQuestion{“iPhone 4s再出问题，有用户反映耳机出现间歇性故障，向外打电话时无法听见声音。据苹果官网显示，关于耳机失灵的讨论已经多达近30页，不少用户表示这次故障能够通过重新开机或者拔掉耳机后恢复，但是这仅是暂时解决问题。”（新闻）}

品种少的好处之一，有问题反应快，解决也快。

\BookQuestion{网友：好观点！}

It's not an opinion. It is a fact.（这不是一个观点，而是一个事实。）（2011-11-07）

我十多年前就明白这个道理，从任天堂那里学的。这些年来很惊奇地发现，很少有人能看懂。（2011-11-08）

\BookQuestion{网友：苹果和Windows的用户体验很不一样，我和朋友们的归纳就是两者都很强大（当然还有UNIX系统也很强大），但是苹果让人用近乎白痴的方式使用强大的功能，而Windows让人用近乎专业人士的方式使用强大的功能（UNIX平台的系统是让人用近乎天才的方式使用强大的功能）。}

你这个形容很形象，这正是苹果厉害的地方。苹果还有不少特别厉害的地方，比如品种单一，所以效率高、质量一致性好、成本低、库存好管理等等。我从做小霸王时就追求品种单一，特别知道单一的好处和难度，这个行业里明白这一点并有意识去做的不多，我们现在也根本做不到这一点。比较一下诺基亚，你就马上能明白品种单一的好处和难度了。

诺基亚需要用很多品种才能做到消费者导向，而苹果用一个品种就做到了，这里面功夫差很多啊。（2011-01-22）

% 来源：EPUB part129.xhtml
\subsection{把用户导向真正做到了极致}

苹果把用户导向真正做到了极致，它绝对是我们制造业学习的榜样。（2011-01-08）

iPhone 4s发售第一天我就跑去排队买了一台，出来时碰上国内某著名公司的一个朋友，然后去他们公司看了一眼，回车上时居然把iPhone的面给摔裂了。今天有空去苹果店换了一台，本来还想给太太也买一台的，发现iPhone 4s已经卖空了。现在要买的人必须晚上9点上网订。第二天去拿，稍微晚点恐怕就要再等第二天了。有意思的是，我先说要买一台，苹果店的人告诉我今天没了，然后我说换一台行不，结果换一台有货。iPhone 4s这个卖法，这个季度的财报要创纪录了。（2011-10-20）

\BookQuestion{网友：为什么要买的需要等一天，要换的却有货呢？是因为要优先保证退换的人可以及时拿到吗？看来我是被“教育”得很好的消费者，如果我要去换，人家告诉我和买的人享受同等待遇而不是要等货源宽松再给我，我就谢天谢地了，如果说我还有优先权，那真是惊喜啊！}

你不明白没关系，反正大多数人或公司也不明白，说明你很正常哦。（2011-10-23）

\BookQuestion{网友：是不是高科技公司也只是满足顾客的消费需求？}

不知道什么是高科技公司。在消费者眼里其实没有这个东西。（2010-10-30）

我也常常想苹果现象。乔布斯真是个绝顶聪明的消费者导向大师和市场营销大师。他们在产品上做的很多取舍是让人叹为观止的，同行其他人没人敢做他们做的一些事。（2010-05-25）

\BookQuestion{网友：苹果5代出来了，配置不好。苹果说他们不拼硬件，那他们拼软件吗？}

体验。（2012-09-16）

\BookQuestion{网友：我觉得乔布斯比较牛的是帮用户创造适用的产品，而不是迎合用户的需求。}

这两者其实没本质区别。乔布斯是把自己当用户，所以别人以为他先知先觉而已。哪一个改变游戏规则的成功者不是这样呢？（2012-07-06）

\BookQuestion{网友：苹果是消费者导向的典范。}

苹果确实特别在意消费者体验。（2013-04-08）

\BookQuestion{网友：亲自处理用户来信，这个非常不难，但在这个公司的这个位置上非常难。}

非常不难？！苹果用户那么多，每天来信恐怕是数以千计吧？各种问题都会到库克那里的。前段时间我的Apple Cash账号出问题，我的助手就是直接发邮件给库克，然后第二天就接到苹果的人的电话了。

\BookQuestion{网友：您不会觉得市长信箱是市长在亲自处理吧？}

你肯定觉得自己很聪明是吧？顺便告诉大家一个数据：库克每天早上大概花3个小时处理用户来信。（2022-09-21）

\BookQuestion{网友：一线用户的反馈真是非常重要，我记得丁磊就说过网易的客服绝对不能外包，苹果这个体量，库克亲自处理每一个用户来信，还是让我有点意外。}

那只是客户反映问题的渠道之一，库克大部分也是直接给相应的部门去处理。这样做非常接地气，比较容易感受到用户的需求。（2022-09-24）

苹果一直聚焦的就是用户体验，所以这样的公司才能出伟大的产品。居然有些人会站在死胡同里说苹果的路会越走越窄，呵呵，有趣。（2013-02-18）

\BookQuestion{网友：常常提到好公司的企业文化是以消费者为导向，但是又常说我们不能太在意市场，太在意华尔街和别人的看法。别人和市场的看法也算消费者的看法吧？}

说不在意的是指应该不在意短期行为或期望。比如不在意华尔街的看法指的是不在意华尔街短期的看法。短期的看法是投票器，长期才是称重器。比如华尔街希望苹果出低价机多占领市场份额，而苹果想的是做最好的产品，所以短期就让华尔街失望了。（2013-03-07）

\BookQuestion{网友：Think different的理念是不是影响了苹果更好地做到消费者导向？}

Think different其实非常了不起，长期而言力量是巨大的。Think different其实是内含消费者导向的，因为那是前提。（2020-11-27）

苹果进入一个产品的时机往往不是因为苹果在择时，而是因为苹果在这之前还没准备好或者说苹果还不满意自己的产品。苹果推出产品的最重要的基点是“我到底能为用户做什么”。很多年前，网上疯传苹果会推苹果电视，我说那么大一台电视，苹果能做什么？结果苹果确实没推。（2023-06-13）

苹果的角度一般是看能不能给消费者带来点什么，而不是生意上能不能成功。根据这个观点，我推断出苹果不会推苹果电视机，前些年我还认为苹果车大概率不会在近些年见到，除非技术条件达到了苹果的要求。（2024-02-05）

苹果果然取消了电动车项目！（2024-02-28）

\BookQuestion{网友：记得一开始传苹果要做电动车的时候，大道说过：“苹果不太可能做电动车。”苹果花了那么多时间、人力、物力研究电动车，最终还是放弃了。苹果真是理性啊。大道说的，“长远看不合适的东西最合适的办法就是现在就停下来！”}

10多年前曾经盛传苹果要推电视机，有业内资深人士甚至告诉我他已经看到样机了。我说，你可能不太了解苹果，样机不说明任何问题。同理，差不多的时期，有人问我iPhone会不会出大屏，我说一定会的。尽管后来还是等了3年。（2024-03-02）

《乔布斯传》最后一章“遗产：无比辉煌的创新天堂”中的一段自述，我反复读了很多遍，字里行间跃然纸上的是一个虽有阴暗面却始终闪亮的灵魂。

吉姆·柯林斯在《基业长青》中曾经总结过伟大企业的其中一项特质是利润之上的追求，乔布斯和他的苹果给予了这项特质以最完美的诠释。（2014-01-31）

\BookQuestion{网友：是不是再好的公司如果没有优秀的管理人，迟早会衰落，就像苹果。可是优秀的管理人员迟早会离开公司，那这个公司是不是就不值得投资了？}

苹果衰落了吗？苹果其实从没有像现在这么强大过！这段话的意思是，没有强大的文化，就不会有伟大的公司。乔布斯就是想建立强大的企业文化并且确实做到了！只有强大的企业文化才能不断吸引到好的人才并留住他们。不能理解这点的人们大概就只能炒一辈子股票了。（2014-02-02）

% 来源：EPUB part130.xhtml
\subsection{库克是乔布斯最伟大的发明之一}

\BookQuestion{网友：您是看好苹果公司的文化，还是看好CEO？}

我曾经一直认为乔布斯是个造钟人，最近突然觉得只有乔布斯的苹果是不可能成为今天这个样子的。今天的苹果，乔布斯的作用已经不是那么大了。就算乔布斯不打算再回来，今天的苹果也会靠惯性向前的。还有就是，当年的苹果董事会为什么会在赶走乔布斯后又能请他回来呢？这可不是容易做到的事。（2011-01-23）

《基业长青》这套书专门说过乔布斯是个“报时人”，所以大概2002、2003年第一次看苹果时（当时市值好像只有大约50亿美元，现金也是差不多这个数），我受这本书的影响，没太认真看就直接跳掉了。还好2011年初突然想明白了，觉得乔布斯实际上是个恰好也会报时的“造钟人”，而且库克是个更好的CEO（更理性）。想通这两点后，就开始下手买了……（2020-01-07）

\BookQuestion{网友：是什么事情让您突然悟到乔布斯还是一个很好的造钟人？}

因为我看见那个钟了。（2012-03-29）

其实，库克是个比乔布斯更好的CEO。库克更理性，同时骨子里对乔布斯的追求非常理解。前段时间见到一个投资界的老前辈时，我说我认为库克其实是个更好的CEO，他老人家说：“其实我也这么认为。”《基业长青》里说乔布斯是个报时人，其实乔布斯同时还是个非常好的造钟人，库克其实就是乔布斯最伟大的发明（发现）之一。（2018-08-07）

\BookQuestion{网友：我很了解乔布斯，但对库克有偏见，可能是我内心拒绝承认苹果的CEO已经变更为库克了。我会好好看近十年的发布会。}

乔布斯选库克是绝对有道理的，这其实也是乔布斯对苹果的最大贡献！对库克的不敬其实是因为对乔布斯的不了解。2011年我就是想通了这个道理才决定开始投资苹果的，看起来我确实做了一件对的事情。我能理解库克能成为好的CEO可能跟我自己的经历有关。我本人不那么懂技术（远不如库克），但不影响我创立和经营一家好公司。我跟库克打过交道，我很喜欢他，觉得他是一个绝对值得信任的人。我认为我还会继续持有苹果很久很久的时间。我相信我看到和理解的东西。（2022-10-19）

\BookQuestion{网友：看了《乔布斯传》，终于明白苹果公司为什么伟大了。无论从战略到细节，乔布斯都100\%全力掌控，战略和执行都到位。}

小说的东西不可全信。对多数人来讲，看这本书只是看了个热闹，很多人会得出很片面的结论，比如没有乔布斯的苹果会不行，等等。有趣的是微软得出的结论和我一样，就是没有乔布斯的苹果更厉害。（2012-01-25）

% 来源：EPUB part131.xhtml
\subsection{维持现金中性}

苹果的策略非常清晰，就是维持现金中性，大致意思就是留够运营需要的资金后，多出来的钱都通过派息或者回购返还给股东。苹果应该是没有实质性债务的，但海外赚的钱回美国税比较高，所以苹果可以用美国海外的钱（或可流动资产）做抵押在美国发债。前几年利息低的时候苹果发了很多长债，然后用这些钱的一部分回购了股份，股东们赚大了哈。不过，海外的钱最后怎么回到公司我还没搞明白过，但看到过某些年政府会开个小口子让钱以比较优惠的汇率回美国。我没看过财报，有时间的人可以去研究，我就简单相信库克说的逻辑就好了。（2024-04-30）

\BookQuestion{网友：如果管理层能用好现金最好是少分红不分红，分红还要交税的，如果管理层无法用好现金，那就应该分红。}

苹果的文化不会乱用现金的，就这么简单。（2012-01-27）

\BookQuestion{网友：您觉得回购股票好还是分红好？}

我觉得分红简单且透明。回购则总是和价格有关，公司需要花时间去判断股价是否够便宜，有点难。（2012-02-24）

\BookQuestion{“苹果每股派息2.65美元，回购100亿美元股票。”（新闻）}

这个花钱法只会越花越多嘛。分红也许该再多些才好，不过由于海外钱回美国要打很多税，这样可能也有道理。总的来讲，开始下毛毛雨了。

分红和增长其实没有必然联系，把这两种联系在一起的逻辑不通。现在的苹果更理性，现金太多没用很可能会用于派息，但回购要麻烦些，因为需要判断股价高低。（2012-03-19）

\BookQuestion{网友：段哥说了，今天的苹果依然便宜。那么苹果的超额现金是否应该优先回购自家股票呢？}

对于管理层而言，回购的问题是需要判断股价是否便宜，但派息则更简单。我估计苹果在留够现金后会有个派息和回购的计划，回购只是在一定情况下实施。如果我是苹果的CEO，就简单宣布以后每年的获利拿40\%出来回购， 40\%派息， 20\%留下。（2012-02-07）

\BookQuestion{网友：巴菲特谈苹果股价低迷时，说我若是库克就回购。}

苹果其实一直在回购，在前面100个亿没回购完以前是不应该出新的回购计划的。另外，这么大的回购计划可能也需要时间去制定，不能有纰漏。长期来讲，股价低对苹果没任何坏处，回购到便宜的股票实际上对长期股东也非常好。举个极端的例子：假设苹果股价在这里不动3年，苹果每年赚450亿，3年后加上现在手里的现金资产近1400亿，手里留400亿周转，再去掉分红300亿，苹果可以投2000亿进去回购，那3年后苹果的市值就只有2000亿左右了。那个时候苹果可能一年可以赚500多亿或者600多亿（如果不是更多的话），很难再继续在2000亿市值上待着吧？（2013-03-05）

\BookQuestion{网友：假如苹果只宣布提高分红，而不大量回购股票（一年至少200亿美元以上），是否可以判断苹果对自己的前景不乐观？}

其实任何股票都只有一个真正的买家，那就是公司自己。

苹果现在拥有远多于运营需要的现金，所以苹果一定会用他们认为合适的办法还给股东，这是人们在投资苹果时应该认定的，不然就是投机。之所以认为这是个回购好机会的理由，因为这个价位以上套住的人很多，股价很可能不会反弹得太快，所以苹果有机会买到好价钱。我如果是苹果，现阶段宁愿贷款也要回购，因为分红都比利息高了。等将来有窗口时再将海外的钱汇回来就行。（2013-04-19）

\BookQuestion{网友：苹果那么多现金，您要是库克您会怎么花？}

如果我是库克，用现金的办法和他一样，就是留够需要的尽量用来回购。这次的回购在2015年底以前结束，相信后面还会有很多类似的回购。最好股价别动，然后别人都把股票卖给苹果了，按这个速度，10年内我就是苹果唯一股东啦，然后我会决定下市的。（2013-04-26）

\BookQuestion{网友：这么好的公司，也许巴菲特正在买入苹果股票，很快会到500美元的。}

我不希望苹果股票涨。最好是几年都不涨，这样回购价钱就可以低一些。（2013-04-28）

去年这时总股数是59.648亿，现在是56.735亿，少了2.9亿多股。这样下去20年后苹果公司就是我的了。

\BookQuestion{网友：再回购下去，段总想要低调就很难了，要进各种财富排行榜了。}

20年就可以了。

\BookQuestion{网友：以后可不能把它给私有化了啊！}

如果目前这个价格能够持续10年的话，想不私有化也难啊。（2015-10-29）

\BookQuestion{网友：苹果刚宣布增大回购，市场就提供低价机会，感觉特配合。}

反正每年苹果都是要宣布回购的。和有些企业宣布不一样的地方是，苹果是来真的。希望苹果的股价继续掉，现金流继续强悍，回购当然也是会继续的，10年8年后我可能就真是苹果大股东了哈？（2016-04-28）

\BookQuestion{网友：查了下资料：苹果公司目前已完成回购1000亿美元计划的30\%，从2012年到2018年，苹果已通过股票回购和派息，向股东返现了2750亿美元，其中股票回购就达到了2000亿美元。再加上这1000亿，回购高达3000亿！我觉得这传递了两个信号：苹果认为自己股价不贵；苹果想更多地回馈股东。}

没完全看出你说的信号，但至少可以看到一个事实，就是如果加上已经返还和回购的部分，现在的苹果市值确实早就到了1万亿了。（2018-11-20）

% 来源：EPUB part132.xhtml
\subsection{对自己懂的公司估值没那么难}

我在2011年买苹果的时候，苹果大概3000亿美元市值（当时股价310/7 = 44），手里有1000亿净现金，那时候利润大概不到200亿。以我对苹果的理解，我认为苹果未来5年左右盈利大概率会涨很多，所以我就猜个500亿（去年595亿）。所以当时想的东西非常简单，用2000亿左右市值买个目前赚接近200亿/年、未来5年左右会赚到500亿/年或以上的公司（而且往后还会继续很好）。如果有这个结论，买苹果不过是个简单算术题，你只要根据你自己的机会成本就可以决定了。但得到这个结论非常不容易，对我来说至少20年功夫吧。能得到这个结论，就叫懂了。不懂则千万千万别碰，我有个球友320/7 = 46买了一些，结果一个回调，310/7就卖了（现在苹果加上分红可能早就超过200了），还跟我讲为什么要卖的道理，从此我不再跟他说投资了。（2019-05-20）

\BookQuestion{网友：任何企业都应该有个合理的估值，还是某些企业根本就无法估值？}

大部分人说的估值都是指市场应该给什么价，这个概念的估值不说也罢。如果你不是这个意义上的估值，你自然会明白的。

其实对自己懂的公司估值没人们说的那么难，难是难在不容易找到自己想买的公司。我对苹果的估值就是个例子，只有两个方面，简单来说就几句话（大概）：

1.3000多亿市值，1000多亿净现金，一年净利润400亿，未来会到500亿或更多（看不懂这个的为什么要对苹果感兴趣？）；

2. 苹果会将利润以合适的方式还给股东。

第二条在A股比较难解，这也是当时博友问我会选苹果还是茅台时，我毫不犹豫选苹果的原因。记得两年前有个朋友在球场碰到我，很沮丧地说他600多买了苹果，那时苹果一路在往下掉。我当时只说了一句：也许两年后你会觉得600很便宜哈。（2014-11-15）

\BookQuestion{网友：以前400-500美元的时候，我是能看出便宜，现在近1000美元了，贵吗？}

现在其实也不算贵，只是没有以前便宜了。就我自己而言，拿着苹果的选择如下：1. 继续持有，享受苹果的成长和成就。现在拿着苹果，在可预见的未来复合回报应该还可以有8\%或以上的（这里指的是企业获利）。2. 找到比苹果更好回报的投资。我目前还没有任何目标，主要是自己能搞懂的生意太少。3. 卖了拿着现金吃利息？这个显然不是太靠谱。当然，有人会说现在卖了等回调再买回来，这么说的人如果有过苹果的话应该早就等得胡子都白了吧？（2015-02-27）

\BookQuestion{网友：您曾经说过苹果的市场份额会到40\%，是因为快消品超过40\%都很罕见吗？这几年下来苹果确实还没到40\% （日本、美国到了）。}

40\%只是毛估估的想法，理由是iPhone虽好，但安卓系统要便宜很多，而且现在的安卓比当年诺基亚的智能机已经好用太多了。这是个多元的世界，总有人会因为各种原因选择不同产品的。我的40\%的意思是，如果苹果能有40\%的份额的话，这个投资就是极好极好的投资了，而且40\%是有可能达到的，因为在很多发达国家或地区目前已经达到40\%或以上了。实际上，苹果只要有25\%以上的份额就足够好了。（2015-04-26）

\BookQuestion{网友：持有容易吗？阿段2011年买了苹果公司后，2012—2013年出现了55\%的下跌；2015—2016年出现了36\%的下跌；2018年出现了40\%的下跌；2020年出现了36\%的下跌。美股，几乎单仓一只。甚至阿段身边的朋友知道阿段还持有苹果的情况下，都忍不住卖掉。}

是啊，苹果每次大跌的时候我都会加码，每一次！想到10年后的时候，自己喜欢的股票大跌怎么会心情不好呢？不过确实绝大部分早年跟我买苹果的都早卖掉了。（2020-11-06）

前些年一直有人在说，树不会长到天上去的。他们是对的，树确实不会长到天上去，但这不会影响苹果的成长。（2021-12-08）

在我的理解里，苹果的商业模式非常强大，未来的复利回报大概率是会高过长期美国国债的。我总是习惯把美国长期国债作为第一机会成本（就是最低的意思，我如果看不到公司的长期盈利是能超过长期国债的，就会直接过滤掉），苹果第二。所以除非看到未来几十年会好过苹果的公司，不然我不愿意换的。（2022-01-27）

\BookQuestion{网友：现在持有现金不急于买入苹果是怎么样的思考角度？}

这个问题巴菲特其实也没完全说清楚过。我是这么想的：我手里的苹果如果卖掉的话，买什么呢？目前看来只能买短期债券。我的现金目前买什么呢？苹果还是短期债券？当然，如果你有其他更好的投资标的就没有这个问题。我今天其实买了10万股苹果，价格是176美元多一点，钱已经付了，但要一年后交货（就是按现价买了苹果同时卖了一个一年后到期的200的看涨期权，收了21美元。最高收益是赚25美元/股）。今天我有很多期权到期，我把赚到的钱的一部分买了苹果，但这并不是这么花钱的理由。我想得最多的是176美元的苹果便宜吗？按这个价钱你愿意买下这个公司吗？其实对阿里我也是在想同样的问题，1800亿你愿意买下这个公司吗？腾讯我想的也是这个问题。很多问题不是马上能够得到结论的，好在我们有时间。需要考虑什么因素？当然是所有因素啦。（2023-12-16）

AI iPhone大概率会导致一波换机潮，主要是因为iPhone太耐用，平均换机周期比较长，很多人都没啥换机的动力。AI iPhone会给大家一个很好的换机理由。不过，换完一轮后大概又会进入相同的换机周期甚至更长的换机周期。估计未来十年内iPhone至少能卖个25亿台或者接近30亿，活跃用户估计也会从13-14亿增加到接近20亿。未来10年里iPhone这个产品大概率会给苹果带来一万亿以上的利润。毛估估乱猜苹果10年后的年利润可能会接近或达到2000亿美元。（2024-07-22）

前段时间我就说过苹果确实不便宜，巴菲特换点到石油是非常容易理解的事情。但目前我也没有太强的意愿拿苹果去换别的公司，主要是我能理解的公司太少。不过，我最近确实花了一点时间去想西方石油（OXY）和雪佛龙这两个石油公司。从占有资源的角度看，这两家公司似乎确实是蛮有价值的。所以我卖了不少苹果的call，同时也卖了不少OXY的put。

一般来说，巴菲特开始减持的公司是会继续减持的。我不会卖苹果，但很乐意看到巴菲特先生减持一些或更多。（2024-02-15）

\BookQuestion{网友：能详细讲讲西方石油的逻辑吗？巴菲特这么长时间频频加仓，可是市场上都是看空原油的。}

其实就是按不错的价钱买了个油田。石油资源大概率是会紧缺的，只是不知道什么时候而已。（2025-03-04）

山也还是那座山，梁也还是那道梁！

今天开始买了一些（或者卖put）苹果、英伟达、谷歌、台积电，一会儿再卖点腾讯的put。

感觉起码要折腾一阵子，卖点put来回一起折腾也许是个好玩的游戏。还是老话，千万别用margin。如果put进来的价格从10年后看回来是便宜的价格，那这个投资就值得做，长短皆宜。（2025-04-08）

\BookQuestion{网友：苹果受关税影响挺大的。}

10年后看也许没那么大。（2025-04-09）
